\section{A cylinder estimate and compactness of the components}
\label{sec:cylinder}

The normalizations of the divisor components have genus zero by
Proposition~\ref{curve:genus}. We prove a cylinder estimate and use it
twice to exclude noncompact normalizations. First, on a planar domain
omitting at least two points of the sphere, the estimate controls a
nonzero scalar differential, contradicting the scalar residue
obstruction. The only remaining open normalization is the plane.
There a periodic gauge lets us apply the same estimate to bundle-valued
differentials on infinitely many translated curves, contradicting
finite-dimensionality of their fixed cohomology space.

The estimate controls the integral of a differential over a specified
target height slab by the curve's area in a fixed enlargement of that
slab. Its constants must remain bounded along arbitrarily long chains
of low-energy source cylinders.

We retain the exponent $A>0$ of Lemma~\ref{curve:area}.
It was fixed from the coordinate geometry before the holomorphic structure
on $E$ was chosen.  Its correction exponent satisfies
$\gamma>2A+10$, as prescribed in \eqref{curve:gamma-choice}.
The normalization of any component of $D$, and of any fixed deck
translate of that component, satisfies
\begin{equation}\label{cylinder:area-input}
 \Area\bigl(B\cap h^{-1}([-\ell-1,-\ell])\bigr)
       \leq C_B e^{A\ell}\qquad(\ell\in\NN).
\end{equation}
The constant $C_B$ may depend on the translate.  The exponent does not.

\subsection{Estimates on source cylinders}

Write $\mathcal C=\RR\times(\RR/2\pi\ZZ)$, with coordinate
$v=t+i\theta$ and flat area $dA_v=dt\,d\theta$.  Put
\[
 S_j=[j,j+1]\times(\RR/2\pi\ZZ),\qquad
 C_j^r=(j-r,j+1+r)\times(\RR/2\pi\ZZ).
\]
Strip boundaries have zero area and can be included when specifying a
point that a strip contains.  For a holomorphic map $f$ into $Y$ let
$u(v)=|df_v(\partial_t)|$.  Conformality gives
\begin{equation}\label{cylinder:energy}
 f^*\omega=u^2dA_v,
\end{equation}
where $\omega$ is the Hermitian area form on complex tangent lines.
Thus $u$ differs from the Hilbert--Schmidt norm $|df|$ only by the fixed
factor $\sqrt2$.  We use $u$ in the estimates to avoid this factor.

A finite nested coordinate cover of $X$, lifted to $Y$ and recentered
inside its smaller members, gives a radius $r>0$ and centered
holomorphic coordinate balls $U_p(4r)$ about every $p\in Y$.
Fix these charts so that the Euclidean and lifted tangent norms are
comparable with one constant $K\geq1$. Decreasing $r$ once if necessary
makes all the recentered balls lie in the larger original charts.
The smaller concentric balls $U_p(r)$ will be used in the definition
of a flat strip.

\begin{lemma}[Small energy fills a source cylinder]
\label{cylinder:small-energy}
There are an integer $P\geq2$ and positive constants
$Q,\delta,\varepsilon_0,D_0$, depending only on the lifted geometry,
with the following property.  Suppose that $\Omega\subset\mathcal C$
is a domain and $f:\Omega\to Y$ is a nonconstant proper holomorphic
map. A strip $S_j$ is called \emph{flat} if
$\overline{C_j^{2P+2}}\subset\Omega$ and the image of this entire
thickened cylinder under $f$ is contained in $U_p(r)$ for one of the
fixed coordinate charts. Every flat
strip satisfies $\sup_{S_j}u\leq D_0$.

For every interval $I\subset\RR$, if $S_j$ contains a point $v_j$
with $h(f(v_j))\in I$ and $S_j$ is not flat, then
\begin{equation}\label{cylinder:energy-cost}
 \int_{\Omega\cap C_j^Q\cap(h\circ f)^{-1}(I+[-\delta,\delta])}
              u^2\,dA_v\geq\varepsilon_0.
\end{equation}
The same assertion holds on a source half-cylinder whenever
$\overline{C_j^Q}$ lies in its interior.
\end{lemma}

\begin{proof}
Choose $\delta$ strictly larger than twice the maximum height
oscillation of the balls $U_p(4r)$, with a fixed positive margin;
a uniform multiple of
$8Kr\|dh\|_\infty+1$ suffices after increasing the chart comparison
constant.

Choose $P\geq2$ large enough that
\begin{equation}\label{cylinder:q}
 q:=\frac{2K^2}{e^P-1}\leq\frac12,
\end{equation}
and set $Q=3P+4$.  Cauchy's estimate on source disks of radius one,
inside $C_j^{2P+2}$, gives a uniform bound $D_0$ for $u$ on
$\overline{C_j^{2P}}$ whenever $S_j$ is flat.  The norm of a coordinate
map with image in $U_p(r)$ is at most $r$, so this bound depends only
on $K,r$, and the fixed disk radius.  In particular it bounds $u$ on
$S_j$.

It remains to prove a uniform positive energy cost. The issue is that
the source domain may omit part of the thickened cylinder. We show
that a sequence with energy tending to zero has a horizontal graph
limit; proper projection near that limit then forces the missing
source points to be present. If no
$\varepsilon_0$ works, there are maps $f_\nu$, nonflat strips, and
marked points as in the statement for which the integral in
\eqref{cylinder:energy-cost} tends to zero. Translate source indices
to $j=0$ and translate the targets and their height intervals by the
same powers of $T$ so that the marked images lie in a fixed compact set.
After a subsequence, the marked source points tend to
$v_\infty\in S_0$ and their images tend to $p\in Y$. Denote the
corresponding translated intervals by $I_\nu$; their lengths play no
role in what follows.
Work in the fixed coordinate ball $U_p(4r)$. For large $\nu$, the marked image is
arbitrarily close to $p$, so the strict margin in $\delta$ makes every
point of this ball have height in the corresponding
$I_\nu+[-\delta,\delta]$, even when the marked height lies at an
endpoint of $I_\nu$. The energy of the graph
portions in
$C_0^Q\times U_p(4r)$ therefore tends to zero in the target direction.

These portions
\[
 \Gamma_\nu=\{(v,f_\nu(v)):v\in\Omega_\nu\cap C_0^Q,
                                  \ f_\nu(v)\in U_p(4r)\}
\]
are closed analytic curves in that product.  To check closedness,
suppose $(v_k,f_\nu(v_k))$ converges in a compact subset of the product.
Properness of $f_\nu$ puts the $v_k$ in a compact subset of
$\Omega_\nu$; their limit is consequently in its domain and belongs
to the graph. On a half-cylinder obtained by cutting a proper full
normalization, the same proof works on compact source disks away from
its artificial finite boundary; properness is applied to the ambient
map. This is the place where properness is essential.

The product Hermitian area of a graph is the sum of its horizontal and
vertical areas.  Its horizontal area is at most $\Area(C_0^Q)$,
because the source projection has at most one point in each fiber.
Its vertical area tends to zero.  Bishop's compactness theorem
\cite{Bishop1964}, in the analytic-cycle and support formulation of
\cite[Theorems~3.9--3.10, pp.~585--586]{HarveyShiffman1974}, therefore
gives a subsequence converging locally as positive analytic cycles,
and with convergence of their supports, to a nonempty pure
one-dimensional analytic cycle $\Gamma_\infty$ through
$(v_\infty,p)$.  These theorems apply in local product coordinates
(equivalently, identify the finite cylinder with an annulus).
The locally bounded graph areas give their mass hypothesis.
Convergence of positive currents makes the limit's vertical area zero.
On the regular part of each component this says that the
target coordinate functions have zero derivative, so each component
is horizontal.  Analytic continuation in the connected cylinder then
gives
\[
 C_0^Q\times\{p\}\subset\supp\Gamma_\infty.
\]
For completeness, horizontal components cannot disappear at an
interior source point: the defining holomorphic functions of the
limit, restricted to the horizontal slice, vanish on an open set and
hence on the whole connected slice.

We spell out why convergence to this slice fills a smaller cylinder.
Fix a source disk $V\Subset C_0^Q$.  Choose a small target ball
$B(p,\eta)\Subset U_p(r)$ whose closure meets no other horizontal
component of the limit.  Such a ball exists by local finiteness of
the analytic cycle.  On a slightly larger source disk the approximating
graphs eventually avoid the target boundary $\partial B(p,\eta)$:
otherwise a convergent sequence of graph points would give a limit
point on that boundary.  Thus the projection of the graph portion
inside $V\times B(p,\eta)$ is proper.  It is finite, since its fibers
have at most one point, and it has nonempty image by convergence to
the horizontal slice.  A nonconstant finite holomorphic map of
analytic curves is open, and a proper map is closed; its image is
therefore all of the connected disk $V$.  Its degree is exactly one,
by the same single-point fiber bound.  In particular every point of
$V$ belongs to $\Omega_\nu$, and its image is in $B(p,\eta)$.
If there were a second horizontal component, the same argument in a
disjoint target ball would give two graph points over one source
point.  Thus there is no second component.

Cover $\overline{C_0^{2P+2}}$ by finitely many such disks.  It follows
that for all sufficiently large $\nu$ this entire compact cylinder
belongs to $\Omega_\nu$ and its image is in $U_p(r)$.  Interior Cauchy
estimates also give convergence of its coordinate derivatives to
zero on smaller cylinders.  In particular $S_0$ is flat, contradicting
its choice.  This proves \eqref{cylinder:energy-cost} and all the
asserted uniformities.
\end{proof}

Nonflat strips therefore cost a fixed amount of area near any marked
height. For the flat strips we use a different bound: along an
arbitrarily long consecutive chain, the sum of the derivative suprema
remains bounded. This will also control the variation of target height.

\begin{lemma}[Flat chains]
\label{cylinder:flat-chains}
Let $J\subset\ZZ$ be a maximal interval of flat indices for a map as
in Lemma~\ref{cylinder:small-energy}.  Put $m_j=\sup_{S_j}u$.
There is a constant $L_0$, depending only on the fixed geometry,
such that
\begin{equation}\label{cylinder:chain-sum}
 \sum_{j\in J}m_j\leq L_0.
\end{equation}
A doubly infinite chain is impossible for a nonconstant map.  If a
chain contains a point with height in an interval $I$, every finite
endpoint of the chain has a neighboring nonflat strip containing a
point with height in $I+[-\Lambda,\Lambda]$, where
$\Lambda=\|dh\|_\infty L_0$ is fixed.
\end{lemma}

\begin{proof}
Take a flat strip $S_j$, and express $f$ in its coordinate chart as
a vector-valued periodic holomorphic function $F(v)$.  Write
$H=\partial_vF$.  Since $F$ is single valued around the cylinder,
\[
 \int_0^{2\pi}H(t+i\theta)\,d\theta
   =\frac1i\int_0^{2\pi}\partial_\theta F(t+i\theta)\,d\theta=0.
\]
Consequently the Laurent--Fourier expansion is
$H(v)=\sum_{n\ne0}c_ne^{nv}$.  Bound the positive coefficients on
the circle $t=j+1+P$ and the negative coefficients on $t=j-P$.
For $j\leq t\leq j+1$ their geometric sums give
\[
 |H(t+i\theta)|\leq
 \frac{2}{e^P-1}
 \max\left\{\sup_{t=j-P}|H|,\sup_{t=j+1+P}|H|\right\}.
\]
The statement is valid for vector values by taking norms in the
Fourier coefficient integrals.  Comparing the coordinate and target
norms yields
\begin{equation}\label{cylinder:contraction}
 m_j\leq q\max_{|k-j|\leq P}m_k
\end{equation}
whenever all the indices on the right belong to the chain.  Each
application uses the chart for its own central strip; the numbers
$m_k$ are intrinsic, so changing charts causes no further loss.

All $m_j$ in the chain are at most $D_0$.  Iterating
\eqref{cylinder:contraction} inside $J$ therefore bounds $m_j$ by
$D_0q^{\lfloor d(j)/P\rfloor}$, where $d(j)$ is the integer distance
to the nearer finite end of $J$.  On a half-infinite interval use its
single finite end.  At most $2P$ indices have any prescribed value of
$\lfloor d(j)/P\rfloor$, so one may take
\[
 L_0=\frac{2PD_0}{1-q}.
\]
If $J=\ZZ$, the same iteration is possible arbitrarily often at every
index and gives $df=0$; analytic continuation contradicts
nonconstancy.

Starting at a marked point of the chain, move in the $t$ direction
at its fixed angular coordinate toward either finite endpoint.
The change of $h\circ f$ is at most
$\|dh\|_\infty\sum_{j\in J}m_j\leq\Lambda$.
The shared boundary with the neighboring nonflat strip is in the
domain because the endpoint flat strip has a full thickening.
That boundary point belongs to the neighboring closed strip and has
height in $I+[-\Lambda,\Lambda]$.  This proves the charging assertion.
Figure~\ref{fig:flat-chain} illustrates the derivative decay along the
chain and the charge to a neighboring nonflat strip.
\end{proof}

\begin{figure}[tb]
\centering
\begin{tikzpicture}[x=0.84cm,y=0.75cm,>=Stealth,
  every node/.style={font=\small}]
 \fill[orange!25] (0,0) rectangle (1,1);
 \fill[blue!9] (1,0) rectangle (10,1);
 \fill[orange!25] (10,0) rectangle (11,1);
 \foreach \x in {0,...,11} {\draw[black!55] (\x,0)--(\x,1);}
 \draw[black!55] (0,0) rectangle (11,1);
 \node[anchor=south,font=\footnotesize] at (0.5,1.08) {nonflat};
 \node[anchor=south,font=\footnotesize] at (10.5,1.08) {nonflat};
 \node[fill=white,inner sep=2pt] at (5.5,0.5)
   {consecutive flat source strips};
 \draw[blue!65!black,thick]
   (1.1,2.75) .. controls (2.1,1.55) and (3.4,1.25) .. (5.5,1.22)
   .. controls (7.6,1.25) and (8.9,1.55) .. (9.9,2.75);
 \node[fill=white,inner sep=2pt] at (5.5,2.35)
   {schematic bound for $m_j=\sup_{S_j}u$};
 \draw[->,blue!65!black] (2.4,1.85)--(3.5,1.42);
 \draw[->,blue!65!black] (8.6,1.85)--(7.5,1.42);
 \draw[<->] (1,-0.35)--(10,-0.35);
 \node[anchor=north] at (5.5,-0.48) {chain length may be arbitrarily large};
 \draw[->] (-0.15,-1.2)--(11.3,-1.2)
   node[anchor=west] {$t$};
 \node[anchor=north,font=\footnotesize] at (5.5,-1.28)
   {source coordinate $v=t+i\theta$};
\end{tikzpicture}
\caption{A flat chain in the source cylinder. The absence of a constant
angular mode makes derivative bounds decay away from the chain's ends.
Their sum, and hence the variation of target height along the chain,
is bounded independently of its length. A chain meeting a target height
slab can therefore be charged to a neighboring nonflat strip in a
fixed enlarged slab. The horizontal coordinate is the source
coordinate $t$, not the target height $h\circ f$.}
\label{fig:flat-chain}
\end{figure}


\begin{proposition}[The cylinder estimate on a target height slab]
\label{cylinder:estimate}
There are constants $R_0,C_0>0$ depending only on the lifted geometry
with the following property.  Let $f:\Omega\to Y$ be a nonconstant
proper holomorphic map from a domain in $\mathcal C$.  Let $w\geq0$
be measurable and suppose
\[
 M_w:=\sup_{j\in\ZZ}\|w\|_{L^2(\Omega\cap S_j)}<\infty.
\]
For every interval $I$ of length one, set
\[
 E_I=\int_{\Omega\cap(h\circ f)^{-1}(I+[-R_0,R_0])}
                          u^2\,dA_v.
\]
If $E_I<\infty$, then
\begin{equation}\label{cylinder:estimate-formula}
 \int_{\Omega\cap(h\circ f)^{-1}(I)}u\,w\,dA_v
            \leq C_0M_w(1+E_I).
\end{equation}
The same estimate holds on a full source half-cylinder cut from a
proper normalization map; properness is used only on thickened strips
away from the artificial initial boundary, and a fixed finite initial
block of strips is included in the constant. The integral on the left is
always restricted to the inverse image of the target interval $I$.
\end{proposition}

\begin{proof}
Take $R_0=\Lambda+\delta$, with constants from the two preceding
lemmas.  Let $\mathcal N_I$ be the set of nonflat strips containing
a point of height in $I+[-\Lambda,\Lambda]$.  The source cylinders
$C_j^Q$ have overlap at most $2Q+3$.  Summing
\eqref{cylinder:energy-cost} over any finite subset of
$\mathcal N_I$ gives
\begin{equation}\label{cylinder:nonflat-count}
 \#\mathcal N_I\leq(2Q+3)\varepsilon_0^{-1}E_I.
\end{equation}
This also proves that the set being counted is finite.  In particular,
the set $\mathcal N'_I$ of nonflat strips hitting $I$ has this bound.
Cauchy--Schwarz, first on each strip restricted to the target slab and
then on the finite sum, gives
\[
 \begin{split}
 &\sum_{j\in\mathcal N'_I}
       \int_{\Omega\cap S_j\cap(h\circ f)^{-1}(I)}u\,w\,dA_v\\
 &\qquad\leq
 M_w(\#\mathcal N'_I)^{1/2}
 \left(\sum_{j\in\mathcal N'_I}
       \int_{\Omega\cap S_j\cap(h\circ f)^{-1}(I)}u^2dA_v\right)^{1/2}
 \leq C M_w E_I.
 \end{split}
\]
Only the intervals appearing explicitly in these integrals are used;
no estimate of the energy outside the enlarged target interval is
being assumed.

On one flat chain, \eqref{cylinder:chain-sum} and
$\|w\|_{L^1(S_j)}\leq\sqrt{2\pi}M_w$ give
\[
 \sum_{j\in J}\int_{S_j}u\,w\,dA_v
      \leq\sqrt{2\pi}M_wL_0.
\]
Every chain meeting $I$ has a finite endpoint adjacent to a member
of $\mathcal N_I$, by Lemma~\ref{cylinder:flat-chains}.
Each nonflat strip has at most two adjacent chains.  Thus at most
$2\#\mathcal N_I$ chains contribute, and their integral over the target
slab is bounded by their full chain integrals. Combining this bound
with \eqref{cylinder:nonflat-count} proves
\eqref{cylinder:estimate-formula} on the full cylinder.

For a half-cylinder, omit the fixed finite block of at most $2Q+4$
strips on which the small-energy neighborhoods meet its initial
boundary.  On this block Cauchy--Schwarz gives
$CM_wE_I^{1/2}\leq CM_w(1+E_I)$.  Apply the preceding argument to the
remaining indices.  There is at most one chain reaching the new
finite endpoint, and its integral is at most
$\sqrt{2\pi}M_wL_0$.  All other chains are charged as before.
Increasing $C_0$ proves the half-cylinder assertion.
\end{proof}

For a proper normalization the integral in \eqref{cylinder:energy}
counts the area of the reduced curve once, away from its discrete
singular set.  A fixed enlargement of a height interval is covered
by a fixed finite number of unit height intervals.  Consequently
\eqref{cylinder:area-input} and Proposition~\ref{cylinder:estimate}
imply, with the same exponent $A$,
\begin{equation}\label{cylinder:exponential}
 \int_{(h\circ f)^{-1}([-\ell-1,-\ell])}u\,w\,dA_v
          \leq C_{B,w}e^{A\ell}\qquad(\ell\in\NN).
\end{equation}
On a half-cylinder the integral is also restricted to that end.
Finitely many height intervals outside the range where the area
estimate was initially proved change only the constant.

\subsection{Excluding nonplane normalizations}

We first apply the estimate with the constant multiplier $w=1$.
This supplies a differential with controlled current mass whenever
the planar normalization omits at least two points of the sphere.

\begin{corollary}\label{cylinder:plane}
If a component of $D$ is noncompact, its normalization is biholomorphic
to $\CC$.
\end{corollary}

\begin{proof}
By Proposition~\ref{curve:genus}, the normalization $R$ is a domain in
$\PP^1$.  If it omits two distinct points, send them by a M\"obius
transformation to $0$ and $\infty$.  Then $R\subset\CC^*$ is a domain
in the logarithmic cylinder and
$\alpha=dz/z=dv$ is a globally defined nonzero holomorphic
differential on $R$.  Formula~\eqref{cylinder:exponential} with
$w=1$ bounds the mass of $f_*\alpha$ on every unit height interval by
$C_Be^{A\ell}$: pulling back a test one-form of norm at most one
costs at most a fixed multiple of $u$, while $|dv|$ is constant.
The scalar case of Proposition~\ref{curve:residue} forbids this
nonzero differential.  Thus $\PP^1\setminus R$ has at most one point.
It has at least one point because $R$ is noncompact.  Removing that
single point gives $R\simeq\CC$.
\end{proof}

\subsection{Excluding plane normalizations}

The scalar residue test leaves only plane normalizations. On a
translated plane the differential $dz$ has an exponentially growing
coefficient in the logarithmic coordinate. We multiply it by the
section of $E$ that tends to zero at the negative end. To control the
product, we first solve the pulled-back Dolbeault equation in a periodic
gauge whose exponential factors have bounded $L^2$ norms on source
strips. The next lemma supplies this gauge, including its angular
constant mode. Its input $a$ is an ordinary complex-valued function;
in the application it will be the coefficient of the pulled-back
connection.

\begin{lemma}[A periodic gauge on the cylinder]
\label{cylinder:gauge}
If $a\in L^1(\mathcal C)\cap L^2(\mathcal C)$, there is a periodic
distribution $g$ satisfying
\[
 \tfrac12(\partial_t+i\partial_\theta)g=a,
 \qquad
 \sup_{j\in\ZZ}\|g\|_{W^{1,2}(S_j)}
       \leq C\bigl(\|a\|_{L^1}+\|a\|_{L^2}\bigr).
\]
For every $p<\infty$ this solution satisfies
\begin{equation}\label{cylinder:gauge-exp}
 \sup_j\int_{S_j}\exp\bigl(p|\operatorname{Re}g|\bigr)dA_v<\infty.
\end{equation}
If $a$ is smooth in an open set, $g$ is smooth there.
\end{lemma}

\begin{proof}
Let $a_0(t)=(2\pi)^{-1}\int_0^{2\pi}a(t,\theta)d\theta$.
Its primitive
\[
 g_0(t)=2\int_{-\infty}^t a_0(s)ds
\]
is bounded by $\pi^{-1}\|a\|_{L^1}$, and
$g_0'=2a_0\in L^2(\RR)$.  This bounds its full $W^{1,2}$ norm on
each unit strip, including its mean.  For $a_\perp=a-a_0$, take
Fourier series in $\theta$ and Fourier transform in $t$.  For every
integer $n\ne0$ define
\[
 \widehat{g_\perp}(\xi,n)
       =\frac{2\widehat{a_\perp}(\xi,n)}{i\xi-n}.
\]
Since $|n|\geq1$, each of the multipliers
$1/(i\xi-n)$, $\xi/(i\xi-n)$, and $n/(i\xi-n)$ is bounded in
absolute value by one.  Plancherel therefore gives
$\|g_\perp\|_{W^{1,2}(\mathcal C)}\leq C\|a\|_{L^2}$.
The sum $g=g_0+g_\perp$ has the asserted periodicity, equation, and
stripwise bound.  Local elliptic regularity for
$\partial_t+i\partial_\theta$ gives the last assertion.

Let $B$ be a uniform bound for these stripwise norms.  The
two-dimensional Trudinger inequality on a fixed bounded strip
\cite{Trudinger1967}, including the $L^2$ term in the norm, gives
constants $c,C>0$ independent of $j$ such that
\[
 \int_{S_j}\exp\left(\frac{c|\operatorname{Re}g|^2}{B^2}\right)dA_v
       \leq C
\]
when $B>0$.  One can cut the cylinder at a fixed angle and apply the
bounded-rectangle form of that inequality; no boundary value of $g$
is prescribed.  The zero-norm case is immediate.  For each finite
$p$, the inequality
$p|x|\leq cx^2/B^2+p^2B^2/(4c)$ proves
\eqref{cylinder:gauge-exp}.  In particular both
$e^{\operatorname{Re}g}$ and $e^{-\operatorname{Re}g}$ have bounded
$L^p$ norms on unit strips, for every fixed finite $p$.
\end{proof}

\begin{theorem}\label{cylinder:compactness}
Every irreducible component of $D$ is compact and has rational
normalization.
\end{theorem}

\begin{proof}
It remains, by Proposition~\ref{curve:genus} and
Corollary~\ref{cylinder:plane}, to exclude a component with proper
normalization $f:\CC\to B\subset Y$.  Let
$M_B=\sup_Bh<\infty$ and $M_D=\sup_Dh<\infty$.
The curves $B_k=T^kB$ are distinct, since their height suprema are
$M_B+k$.  For every sufficiently large positive integer $k$ this
supremum exceeds $M_D$, so $B_k$ is not contained in $D$.
Write $f_k=T^k\circ f$.  We will construct on each of these curves
the nonzero $f_k^*E$-valued holomorphic differential
\begin{equation}\label{cylinder:bundle-differential}
 \alpha_k=(f_k^*s)\,dz
\end{equation}
and prove that its pushforward has growth exponent $A$.

The key point is the cancellation of the growth of $dz$. Properness
sends the exterior of the plane to the negative end of $Y$. After a
periodic change of frame there, the section coefficient becomes an
ordinary holomorphic function tending to zero at infinity, hence of
size $O(1/|z|)$. Since $dz=z\,dv$, this leaves only an exponential
gauge factor, whose stripwise $L^2$ bound is exactly the multiplier
hypothesis in Proposition~\ref{cylinder:estimate}.

Fix one such $k$ temporarily.  Properness, together with the upper
height bound, implies
\begin{equation}\label{cylinder:proper-end}
 h(f_k(z))\longrightarrow-\infty
                  \quad\hbox{uniformly as }|z|\longrightarrow\infty.
\end{equation}
Indeed, otherwise a sequence with $|z|\to\infty$ would have image in
some compact height slab $\{c\leq h\leq M_B+k\}$, contradicting
properness.  On an exterior disk $|z|>R_k$ we may therefore use the
prescribed smooth negative-end frame of $E$.  Choose its Hermitian
norm to be one on that end.  In this frame write
$\dbar_E=\dbar+\eta$, where
$|\eta|\leq C e^{-\gamma|h|}$; this is the choice
\eqref{curve:gamma-choice}.  In the logarithmic coordinate $v=\log z$,
write $f_k^*\eta=a(v)d\bar v$.  Then
\begin{equation}\label{cylinder:connection-bound}
 |a(v)|\leq C e^{-\gamma|h(f_k(v))|}\,u(v).
\end{equation}

After increasing $R_k$, the entire image of this exterior disk has
negative height.  Summing \eqref{cylinder:connection-bound} over its
target height intervals, using \eqref{cylinder:exponential} with
$w=1$ for the first estimate and \eqref{cylinder:energy} for the
second, gives
\begin{align*}
 \int |a|\,dA_v
   &\leq C_k\sum_{\ell\geq0}e^{-(\gamma-A)\ell}<\infty,\\
 \int |a|^2\,dA_v
   &\leq C_k\sum_{\ell\geq0}e^{-(2\gamma-A)\ell}<\infty.
\end{align*}
An initial compact annulus only adds a finite quantity.  Smoothly
cut $a$ off toward the finite side and extend by zero to the full
cylinder.  This retains both bounds and agrees with $a$ on a
slightly smaller exterior disk.  Lemma~\ref{cylinder:gauge} gives
a single-valued periodic smooth gauge $g$ there, with
$\dbar g=a\,d\bar v$ and uniformly bounded stripwise
$W^{1,2}$ norms.

Let $\sigma(v)$ be the coefficient of $f_k^*s$ in the same smooth
frame.  Holomorphicity says $\dbar\sigma+a\sigma\,d\bar v=0$.
Hence
\[
 H(v)=e^{g(v)}\sigma(v)
\]
is ordinary holomorphic and periodic on the exterior cylinder.
By Proposition~\ref{compact:section} and \eqref{cylinder:proper-end},
$\sup_{S_j}|\sigma|\to0$ as $j\to\infty$.  The exponential
estimate \eqref{cylinder:gauge-exp} therefore yields
\[
 \int_{S_j}|H|^2dA_v
 \leq\left(\sup_{S_j}|\sigma|^2\right)
                    \int_{S_j}e^{2\operatorname{Re}g}dA_v
       \longrightarrow0.
\]
The mean-value inequality on fixed small disks in the cylinder,
each contained in the union of three adjacent strips, implies
$H(v)\to0$ uniformly as $t\to\infty$.  Viewed as a holomorphic
function of $z$, it consequently extends holomorphically across
$z=\infty$ with value zero.  Its Taylor expansion in $1/z$ gives
\begin{equation}\label{cylinder:H-decay}
 |H(z)|\leq C_k|z|^{-1}
\end{equation}
on a further exterior disk.

Since $dz=e^v\,dv$, Equations \eqref{cylinder:bundle-differential}
and \eqref{cylinder:H-decay} give the coefficient estimate
\[
 |\alpha_k|_v=|e^v\sigma(v)|
        \leq C_k e^{-\operatorname{Re}g(v)}.
\]
Here $|\cdot|_v$ is the norm of its coefficient relative to $dv$
and the unit smooth frame.  The right-hand factor has uniformly
bounded $L^2$ norms on unit strips by
\eqref{cylinder:gauge-exp}.  Apply
\eqref{cylinder:exponential} with
$w=e^{-\operatorname{Re}g}$.  Pullback of a test one-form costs at
most $Cu$, so this proves
\begin{equation}\label{cylinder:current-growth}
 \Mass\left((f_k)_*\alpha_k\bigm|_{h^{-1}([-\ell-1,-\ell])}\right)
          \leq C_k'e^{A\ell}\qquad(\ell\geq0).
\end{equation}
The omitted compact disk in the source has compact image and
contributes only on finitely many height intervals, which are
absorbed in $C_k'$.  The section $f_k^*s$ is nonzero because
$B_k\not\subset D$, and $dz$ is nowhere zero; thus $\alpha_k$ is a
nonzero holomorphic bundle-valued differential.

Finally fix, once and for all, an admissible $\eps>0$ and a number
$b>\max\{A,\eps\}$.  The currents in
\eqref{cylinder:current-growth} all belong to the same degree-one
space $W^{-4}_{-\eps,b}$ of the shifted complex from
Theorem~\ref{weighted:currents}; their coefficients take values
in the same bundle $E$.  Their positive supports are bounded
separately, and any finite collection has a common upper support
bound.  Proposition~\ref{curve:residue} therefore makes the classes
of every finite collection linearly independent in the single
space $H^{2,1}_{-\eps,b}(Y,E)$.  That space is finite dimensional
by Theorem~\ref{weighted:currents}.  Infinitely many distinct
$B_k$ contradict this finite dimension.  The constants $C_k'$ may
vary; no bound on their supremum has been used.  The exponent $A$,
the weights, and the Sobolev orders were fixed throughout.
This excludes the plane case and completes the proof.
\end{proof}
